\section{Del Nodo al Módulo}

\subsection{Instanciación de módulos Python puros}

Una vez recuperados los valores del Parameter Server, el nodo los usa para construir los objetos que implementan la lógica de control. En este punto el código abandona por completo el dominio de ROS~2. La Terminal~\ref{lst:pva_inst} muestra cómo se instancia el módulo PVA:

\begin{figure}[H]
\centering
\begin{lstlisting}[style=python, caption={Instanciación del módulo PVA desde el nodo.}, label={lst:pva_inst}]
from control_nodes.algorithms.pva import PVA

self.pva = PVA(
    n_rays_used=n_rays,   # int 5
    v_max=vmax,           # float 0.8
    d_safe=self.get_parameter('d_safe').value,
    d_influence=self.get_parameter('d_influence').value,
    xi=self.get_parameter('xi').value,
    rp=0.25,
)
\end{lstlisting}
\end{figure}

\texttt{n\_rays} es un entero Python normal. La clase \texttt{PVA} no sabe nada de ROS~2, no sabe que existe un launch file, no sabe que hubo una terminal. Recibe un argumento de constructor como cualquier clase Python.

\subsection{La clase PVA como Python puro}

La Terminal~\ref{lst:pva_class} muestra el constructor de la clase \texttt{PVA}, que es Python estándar sin ninguna dependencia de ROS~2:

\begin{figure}[H]
\centering
\begin{lstlisting}[style=python, caption={Constructor de la clase \texttt{PVA}.}, label={lst:pva_class}]
class PVA:
    def __init__(self,
                 d_safe=0.3,
                 d_influence=1.0,
                 xi=1.0,
                 rp=0.25,
                 v_max=0.5,
                 w_max=0.5,
                 n_rays_used=None):

        self.d_safe      = d_safe
        self.d_influence = d_influence
        self.xi          = xi
        self.rp          = rp
        self.v_max       = v_max
        self.w_max       = w_max
        self.n_rays_used = n_rays_used
\end{lstlisting}
\end{figure}

Desde aquí en adelante es orientación a objetos Python estándar. La separación entre el código de ROS~2 (el nodo) y la lógica de la aplicación (el módulo) es una práctica arquitectónica deliberada: permite probar la clase \texttt{PVA} de forma independiente, sin necesitar un contexto ROS~2 activo.

\subsection{Por qué esta separación es importante}

El nodo actúa como adaptador: su responsabilidad es recibir datos de la infraestructura ROS~2 (tópicos, parámetros) y entregarlos a los módulos de lógica de control en los tipos correctos. Los módulos implementan los algoritmos sin acoplarse al middleware. Esto tiene consecuencias concretas. En cuanto a testabilidad, los módulos pueden probarse con \texttt{unittest} sin inicializar \texttt{rclpy}. En cuanto a portabilidad, el mismo módulo puede usarse en otro sistema de middleware cambiando solo el nodo adaptador. En cuanto a claridad, la complejidad de ROS~2 queda confinada al nodo y la lógica de control es matemáticamente pura.
